\chapter{Базисно-укоренённый примитивный цикл полного колчана}
\label{ch:v7-baseline-rooted-primitive-cycle-admission-gate}

\section{Вопрос после рангового запрета}

Ранга-один когерентность не выбрала целевой граф, потому что её опора
прямоугольна. Следующий минимальный инвариант должен быть чувствителен не
только к набору рёбер, но и к порядку их прохождения. Проверяем целевой
шестикромочный цикл
\begin{equation}
 \mathcal C_*=Q_L\to u_R\to X_L\to e_R\to L_L\to Y_R\to Q_L.
 \label{eq:v7-rooted-cycle-target}
\end{equation}
Два его ребра,
\begin{equation}
 e_q=Q_Lu_R,\qquad e_\ell=L_Le_R,
 \label{eq:v7-rooted-cycle-baseline-edges}
\end{equation}
уже принадлежат исходному $H_{15}$. Остальные четыре являются новыми.
Это позволяет спросить, выделяет ли старый физический фон цикл без
проекторов на метки ``старая'' и ``новая'' копия.

\section{Обычный шестой след не является селектором}

Пусть $A(x)$ --- симметричная взвешенная матрица смежности полного графа,
где каждому неориентированному ребру $e$ сопоставлена переменная $x_e$.
Тогда
\begin{equation}
 \Tr A(x)^6
 =\sum_{v_0\sim\cdots\sim v_5\sim v_0}
   \prod_{j=0}^{5}x_{\{v_j,v_{j+1}\}},
 \qquad v_6=v_0.
 \label{eq:v7-rooted-cycle-ordinary-trace}
\end{equation}
След суммирует все замкнутые обходы, включая немедленные возвраты.
Точное разложение содержит $305$ различных мономов. Только $14$ из них
квадратсвободны и соответствуют простым шестикромочным циклам; каждый
имеет одинаковую кратность $12$. Поэтому коэффициент целевого слова не
отличает его от конкурентов.

Этот результат согласуется с общей формулой спектрального действия на
колчанах как суммы по замкнутым путям \cite{v7-rooted-cycle-quiver-action}.

\section{Небэктрекинговая очистка}

Обозначим через $\vec E$ множество ориентированных рёбер. Взвешенный
оператор Хашимото зададим формулой
\begin{equation}
 H(x)_{(a,b),(c,d)}
 =\delta_{bc}(1-\delta_{ad})x_{\{c,d\}}.
 \label{eq:v7-rooted-cycle-hashimoto}
\end{equation}
Он запрещает немедленно проходить то же ребро назад. Связь такого оператора
с дзета-функцией Ихары и взвешенными небэктрекинговыми обходами стандартна
\cite{v7-rooted-cycle-hashimoto-source,v7-rooted-cycle-weighted-ihara}.

На текущем двудольном графе шестой небэктрекинговый след очищается до
\begin{equation}
 \Tr H(x)^6=12\sum_{\mathcal C\in\mathfrak C_6}
 \prod_{e\in\mathcal C}x_e,
 \qquad |\mathfrak C_6|=14.
 \label{eq:v7-rooted-cycle-nonbacktracking-trace}
\end{equation}
То есть возвратные слова исчезают полностью, но четырнадцать простых циклов
по-прежнему имеют одинаковый вес. Небэктрекинг является необходимой
очисткой, но не абсолютным селектором.

\section{Относительное укоренение на старом фоне}

Введём оператор извлечения тех мономов, которые содержат оба исходных ребра:
\begin{equation}
 \mathcal R_6(x)
 =\left(x_{e_q}\frac{\partial}{\partial x_{e_q}}\right)
  \left(x_{e_\ell}\frac{\partial}{\partial x_{e_\ell}}\right)
  \Tr H(x)^6.
 \label{eq:v7-rooted-cycle-response}
\end{equation}
Полный конечный перебор даёт точное тождество
\begin{equation}
 \mathcal R_6(x)
 =12x_{Q_Lu_R}x_{u_RX_L}x_{X_Le_R}
 x_{e_RL_L}x_{L_LY_R}x_{Y_RQ_L}.
 \label{eq:v7-rooted-cycle-unique-word}
\end{equation}
Следовательно, среди всех примитивных шестикромочных циклов ровно один
проходит одновременно через $e_q$ и $e_\ell$. Он совпадает с
$\mathcal C_*$ и не содержит ни одного из пяти нежелательных новых рёбер.

Каноничность здесь относительна. Полный немеченый граф сохраняет четыре
типосохраняющих автоморфизма, порождённых обменами
$L_L\leftrightarrow Y_L$ и $e_R\leftrightarrow X_R$. После фиксации
опоры исходного $D_{H_{15}}$ стабилизатор становится тривиальным:
\begin{equation}
 |\operatorname{Aut}_{\rm type}(G)|=4,
 \qquad
 |\operatorname{Stab}_{\rm type}(G,E_{H_{15}})|=1.
 \label{eq:v7-rooted-cycle-automorphisms}
\end{equation}
Значит, селектор не выводится из полного графа в пустоте, но выводится
относительно физического фона, существовавшего до расширения. Проекторы
старой и новой копии отдельно не вводятся.

\section{Почему это ещё не родитель действия}

Цикл содержит только четыре из шести желаемых новых рёбер. Две
вектороподобные массы
\begin{equation}
 X_LX_R,\qquad Y_LY_R
 \label{eq:v7-rooted-cycle-missing-vector-masses}
\end{equation}
в $\mathcal R_6$ отсутствуют. Нежелательное ребро $X_LY_R$ является
хордой $\mathcal C_*$, но также не получает положительной массы одним
извлечением цикла.

После фиксации ненулевых фоновых амплитуд $e_q,e_\ell$ выражение
\eqref{eq:v7-rooted-cycle-unique-word} имеет степень четыре по новым полям.
Поэтому
\begin{equation}
 \Hess_{z=0}\mathcal R_6=0.
 \label{eq:v7-rooted-cycle-zero-hessian}
\end{equation}
Оно не запускает расширение из нуля, не создаёт две массы и не подавляет
остальные разрешённые рёбра. Добавление этих квадратичных знаков вручную
снова означало бы запись ответа в потенциал.

\begin{theorem}[Относительный допуск примитивного цикла]
На полном четырёхвершинном расширении обычный шестой след и его
небэктрекинговая очистка не выделяют целевой цикл: существуют четырнадцать
простых шестикромочных слов одинаковой кратности. Однако совместное
укоренение на двух рёбрах исходного $H_{15}$ оставляет ровно один цикл,
совпадающий с требуемым смешанным маршрутом. Селектор каноничен относительно
предшествующего фона, но не является полным родительским действием.
\end{theorem}

\begin{nogocriterion}[Граница относительного успеха]
Нельзя объявлять \eqref{eq:v7-rooted-cycle-unique-word} полным селектором
шестирёберного ремонта. Следующий гейт обязан одним производным оператором
получить квадратичный запуск четырёх циклических рёбер, две
вектороподобные массы и положительную щель пяти нежелательных рёбер. Свободные
коэффициенты, ручные проекторы копий и предварительное зануление хорд
запрещены.
\end{nogocriterion}

\begin{protocol}
\begin{enumerate}
 \item вершин полного графа: \textbf{$9$};
 \item рёбер полного графа: \textbf{$14$};
 \item простых циклов длины $4,6,8$: \textbf{$12,14,8$};
 \item бесхордовых шестикромочных циклов: \textbf{$0$};
 \item мономов в $\Tr A^6$: \textbf{$305$};
 \item простых шестикромочных мономов: \textbf{$14$};
 \item их общая кратность: \textbf{$12$};
 \item небэктрекинговых мономов шестой степени: \textbf{$14$};
 \item циклов, содержащих оба корневых ребра $H_{15}$: \textbf{$1$};
 \item нежелательных новых рёбер в нём: \textbf{$0$};
 \item порядок типосохраняющей группы до/после фона: \textbf{$4/1$};
 \item новых циклических рёбер выбрано: \textbf{$4$};
 \item вектороподобных массовых рёбер выбрано: \textbf{$0$ из $2$};
 \item квадратичный запуск из нуля: \textbf{нет};
 \item итог: \textbf{положительный относительный наблюдаемый, не полный
 родитель}.
\end{enumerate}
\end{protocol}

Машинный сертификат сохранён в
\path{s2t_v7_baseline_rooted_primitive_cycle_admission_gate_results.json}.

\begin{thebibliography}{9}
\bibitem{v7-rooted-cycle-quiver-action}
C.~I.~P\'erez-S\'anchez,
\emph{The Spectral Action on Quivers}, arXiv:2401.03705.
\bibitem{v7-rooted-cycle-hashimoto-source}
K.~Hashimoto,
\emph{Zeta Functions of Finite Graphs and Representations of
$p$-Adic Groups}, Advanced Studies in Pure Mathematics 15 (1989),
211--280.
\bibitem{v7-rooted-cycle-weighted-ihara}
M.~Kempton,
\emph{Non-backtracking random walks and a weighted Ihara's theorem},
Open Journal of Discrete Mathematics 6 (2016), 207--226;
arXiv:1603.05553.
\end{thebibliography}